% Propietario conservado: /Users/ruben/Documents/ChatGPT/jueces y controles/REVISION_ARGUMENTAL_APERTURAS_CIERRES_20260915/delivery/ARTICLE_V_ES_ARGUMENTAL_20260915/manuscrito/sections/50_estadisticas_ocupacion.tex
% Recuperación documental para el artículo V; RESULTADO_RECUPERADO.
% Propietario: /Users/ruben/Documents/New project/output/REVISION_EDITORIAL_HMT_MD_20260905/01_FUENTES/INTEGRAL/tree/New project/output/ENSAMBLADO_CAUSAL_RESTAURADO_HMT_MD_20260905/fuente/manuscrito/integracion_83/parte_iv/body/B014_ch55_estadisticas_ocupacion_7549b83967f5_04d_ocupacion_estadistica.tex
% Líneas de origen: 1--225.
% Sólo se adaptan las etiquetas de referencia y, cuando procede, el nivel de títulos.
% La matriz de recuperación conserva la correspondencia y las dependencias.
\section{Enumeración microcanónica de estados sobre una pantalla finita}
\label{v:v:rec:ocupacion:section:01}

Consideremos la realización finita
\[
 \mathcal H_{\rm occ}=\bigoplus_{k=0}^{3}\CC^{g_k},
 \qquad
 H_{\rm occ}|_{\CC^{g_k}}=E_k^{\rm occ} I,
\]
con niveles y multiplicidades
\[
 E^{\rm occ}=(0,3,6,9),\qquad g=(1,6,12,8),
 \qquad(N,E)=(12,66).
\]
La suma \(1+6+12+8=27\) hace de \(\mathcal H_{\rm occ}\) una pantalla de
veintisiete modos. Este espectro es el dato de una representación finita
declarada; su cardinal no basta para identificarla con cualquier otra
aparición de \(27\) en HMT\@. El morfismo que relaciona una pantalla concreta
de frontera con \((\mathcal H_{\rm occ},H_{\rm occ})\) debe especificar tanto
las fibras como el operador espectral.
Las funciones generadoras fermiónica y bosónica son
\[
 F(y,x)=\prod_k(1+yx^{E_k^{\rm occ}})^{g_k},
 \qquad
 B(y,x)=\prod_k(1-yx^{E_k^{\rm occ}})^{-g_k}.
\]

\begin{proposition}[Cardinalidades microcanónicas]
\label{v:v:rec:prop:cardinalidades-microcanonicas}
\[
 [y^{12}x^{66}]F=2\,104\,494,
 \qquad
 [y^{12}x^{66}]B=259\,436\,430.
\]
\end{proposition}
\begin{proof}
En el producto fermiónico, elegir \(r\) de los \(g_k\) modos contribuye
\(\binom{g_k}{r}y^rx^{rE_k^{\rm occ}}\). En el bosónico, distribuir \(r\)
ocupaciones indistinguibles entre \(g_k\) modos contribuye
\(\binom{g_k+r-1}{r}y^rx^{rE_k^{\rm occ}}\). La convolución finita sobre los
cuatro niveles produce los enteros indicados.
\end{proof}

\section{Distribuciones de Bose--Einstein y Fermi--Dirac}
\label{v:v:rec:ocupacion:section:02}

La formulación de equilibrio parte del Hamiltoniano y del número total de
ocupaciones,
\begin{equation}
 H=\sum_kE_k^{\rm occ} n_k,
 \qquad
 N=\sum_k n_k,
 \qquad
 K_\mu:=H-\mu N.
 \label{v:v:rec:eq:hamiltoniano-gran-canonico-rev3}
\end{equation}
Para \(\beta>0\), el estado gran canónico es
\begin{equation}
 \rho_{\beta,\mu}
 =\frac{e^{-\beta K_\mu}}{\mathcal Z_{\beta,\mu}},
 \qquad
 \mathcal Z_{\beta,\mu}
 =\operatorname{Tr}(e^{-\beta K_\mu}).
 \label{v:v:rec:eq:estado-gran-canonico-rev3}
\end{equation}

\begin{theorem}[Factorización gran canónica y leyes de ocupación]
\label{v:v:rec:thm:factorizacion-gran-canonica-rev3}
Para niveles de energía \(E_k^{\rm occ}\) con degeneraciones \(g_k\), la
traza fermiónica factoriza como
\begin{equation}
 \mathcal Z_{\rm FD}
 =\prod_k\bigl(1+e^{-\beta(E_k^{\rm occ}-\mu)}\bigr)^{g_k},
 \qquad
 \langle n_k\rangle_{\rm FD}
 =\frac{g_k}{e^{\beta(E_k^{\rm occ}-\mu)}+1}.
 \label{v:v:rec:eq:factorizacion-fd-rev3}
\end{equation}
En el sector bosónico, bajo la condición
\(\mu<\inf_kE_k^{\rm occ}\),
\begin{equation}
 \mathcal Z_{\rm BE}
 =\prod_k\bigl(1-e^{-\beta(E_k^{\rm occ}-\mu)}\bigr)^{-g_k},
 \qquad
 \langle n_k\rangle_{\rm BE}
 =\frac{g_k}{e^{\beta(E_k^{\rm occ}-\mu)}-1}.
 \label{v:v:rec:eq:factorizacion-be-rev3}
\end{equation}
\end{theorem}

\begin{proof}
El operador \(K_\mu\) es una suma de operadores de ocupación conmutantes, de
modo que su traza factoriza por modos. En cada modo fermiónico se suman las
ocupaciones \(0\) y \(1\), y el factor correspondiente es
\(1+u_k\), donde
\(u_k=e^{-\beta(E_k^{\rm occ}-\mu)}\). En cada modo bosónico se suman todas
las ocupaciones no negativas y se obtiene la serie geométrica
\((1-u_k)^{-1}\); la condición sobre \(\mu\) garantiza \(u_k<1\). Las
potencias \(g_k\) incorporan la degeneración. Finalmente,
\[
 \langle n_k\rangle
 =-\frac1\beta\frac{\partial}{\partial E_k^{\rm occ}}
   \log\mathcal Z
\]
produce las dos leyes de ocupación
\cite{v:Bose1924,v:Einstein1924,v:Fermi1926,v:Dirac1926,v:Pathria2011}.
\end{proof}

La derivación microcanónica recupera el mismo resultado en el régimen
termodinámico. Para ocupaciones fermiónicas medias
\(p_k=m_k/g_k\), la entropía combinatoria por nivel es
\begin{equation}
 S_k\simeq
 g_k\bigl[-p_k\log p_k-(1-p_k)\log(1-p_k)\bigr].
 \label{v:v:rec:eq:entropia-fd-nivel-rev3}
\end{equation}
La variación de \(\sum_kS_k\), sujeta a
\(\sum_kg_kp_k=N\) y
\(\sum_kg_kE_k^{\rm occ}p_k=E\), da
\begin{equation}
 \log\frac{1-p_k}{p_k}=\lambda_N+\beta E_k^{\rm occ},
 \qquad
 \mu=-\frac{\lambda_N}{\beta},
 \label{v:v:rec:eq:variacion-fd-rev3}
\end{equation}
y, por tanto, la ocupación de Fermi--Dirac de
\eqref{v:v:rec:eq:factorizacion-fd-rev3}. Los multiplicadores \(\beta\) y \(\mu\)
son las coordenadas duales de las restricciones de energía y población.
El símbolo \(\lambda_N=-\beta\mu\) designa el multiplicador asociado al
número de ocupaciones; \(\alpha\) conserva en todo el artículo su uso
para la constante de estructura fina.

Para ocupaciones bosónicas medias
\(\overline n_k=m_k/g_k\), la logmultiplicidad del nivel es
\begin{equation}
 S_k\simeq
 g_k\bigl[(1+\overline n_k)\log(1+\overline n_k)
          -\overline n_k\log\overline n_k\bigr].
 \label{v:v:rec:eq:entropia-be-nivel-rev3}
\end{equation}
La variación de \(\sum_kS_k\), sujeta a
\(\sum_kg_k\overline n_k=N\) y
\(\sum_kg_kE_k^{\rm occ}\overline n_k=E\), produce
\begin{equation}
 \log\frac{1+\overline n_k}{\overline n_k}
 =\lambda_N+\beta E_k^{\rm occ},
 \qquad
 \overline n_k
 =\frac{1}{e^{\beta(E_k^{\rm occ}-\mu)}-1},
 \qquad
 \mu=-\frac{\lambda_N}{\beta}.
 \label{v:v:rec:eq:variacion-be-rev3}
\end{equation}
Ésta es la ocupación de Bose--Einstein de
\eqref{v:v:rec:eq:factorizacion-be-rev3}. La condición
\(\mu<\inf_kE_k^{\rm occ}\) asegura la positividad del denominador y la
convergencia de cada factor modal.

En la instancia finita construida, las dos restricciones determinan
\[
 (\beta_{\rm FD},\mu_{\rm FD})
 =(0.15307011353415928,\,4.502865197597043),
\]
\[
 (\beta_{\rm BE},\mu_{\rm BE})
 =(0.05456727960068316,\,-15.92569781906737).
\]
Si \(m_k^{\rm mic}\) denota la ocupación media obtenida por enumeración
microcanónica exacta y \(m_k^{\rm eq}\) la evaluación de la distribución
correspondiente, los errores
\(\sum_k|m_k^{\rm mic}-m_k^{\rm eq}|\) son
\[
 0.09897380039002812\quad\hbox{y}\quad
 0.9190625715582017,
\]
respectivamente. La enumeración recorre todas las cuádruplas de ocupación,
calcula sus multiplicidades enteras y resuelve las dos ecuaciones de
restricción. Las dos leyes se realizan sobre una misma pantalla. La elección
física entre ellas exige una graduación y un Hamiltoniano compatibles con la
representación de la trayectoria. El parámetro térmico es
\(\beta=(k_B\Theta)^{-1}\), con la transducción dimensional desarrollada en
la sección~\ref{v:v:ant:sec:ii-kb-aut-desarrollo}.

\clearpage
\section{Convergencia cuantificada hacia la distribución de Maxwell--Boltzmann}
\label{v:v:rec:ocupacion:section:03}

La relación entre las dos estadísticas cuánticas y el régimen clásico puede
formularse sin recurrir a una aproximación informal. Para cada modo, sea
\[
 u_k:=e^{-\beta(E_k^{\rm occ}-\mu)}.
\]
Las ocupaciones por estado se escriben
\[
 \overline n_k^{\rm FD}=\frac{u_k}{1+u_k},
 \qquad
 \overline n_k^{\rm BE}=\frac{u_k}{1-u_k};
\]
las multiplicidades \(g_k\) se incorporan después mediante
\(\langle n_k\rangle=g_k\overline n_k\).

\begin{theorem}[Límite clásico común de las estadísticas cuánticas]
\label{v:v:rec:thm:limite-clasico-fd-be-rev2}
Si \(0\leq u_k\leq\delta<1\), entonces
\[
 \left|\overline n_k^{\rm FD}-u_k\right|\leq u_k^2,
 \qquad
 \left|\overline n_k^{\rm BE}-u_k\right|
 \leq\frac{u_k^2}{1-\delta}.
\]
En consecuencia, en el régimen diluido \(u_k\to0\),
\[
 \boxed{
 \overline n_k^{\rm FD}\sim
 \overline n_k^{\rm BE}\sim
 e^{-\beta(E_k^{\rm occ}-\mu)}} ,
\]
que es la distribución de Maxwell--Boltzmann por modo.
\end{theorem}

\begin{proof}
Las identidades exactas
\[
 u_k-\frac{u_k}{1+u_k}=\frac{u_k^2}{1+u_k},
 \qquad
 \frac{u_k}{1-u_k}-u_k=\frac{u_k^2}{1-u_k}
\]
proporcionan las cotas. Para \(0<u_k\leq\delta\), al dividir por \(u_k\) y
tomar \(u_k\downarrow0\) se obtiene
la equivalencia asintótica. El primer término de orden \(u_k^2\) conserva la
diferencia física: la exclusión disminuye la ocupación fermiónica, mientras
la acumulación aumenta la ocupación bosónica.
\end{proof}
